% Generated by nanotyrannus/tools/make_cards.js from js/facts.js. Edit the facts there, not here.
\documentclass[11pt]{article}
\usepackage[T1]{fontenc}
\usepackage[default]{sourcesanspro}
\usepackage[letterpaper,top=0.5in,bottom=0.4in,left=0.25in,right=0.25in,nohead,nofoot]{geometry}
\usepackage[table]{xcolor}
\usepackage{array}

\newlength{\cardw}\setlength{\cardw}{3.95in}
\newlength{\cardh}\setlength{\cardh}{2.4in}
\definecolor{cut}{gray}{0.62}
\definecolor{soft}{gray}{0.35}
\arrayrulecolor{cut}
\setlength{\arrayrulewidth}{0.4pt}
\setlength{\tabcolsep}{0pt}
\pagestyle{empty}
\setlength{\parindent}{0pt}

% number, date, text
\newcommand{\factcard}[3]{%
  \parbox[t][\cardh][t]{\cardw}{%
    \vspace{0.13in}\hspace{0.18in}%
    \parbox[t]{\dimexpr\cardw-0.36in}{%
      {\fontsize{34}{34}\selectfont\bfseries #1}\hfill{\fontsize{20}{20}\selectfont\bfseries #2}\par
      \vspace{0.07in}%
      {\raggedright\fontsize{15}{19}\selectfont #3\par}%
    }%
  }%
}

\begin{document}
\centering
\begin{tabular}{|@{}c@{}|@{}c@{}|}\hline
\factcard{1}{1942}{A theropod skull was found in Hell Creek Formation rocks from Montana. This formation is known for producing Tyrannosaurus rex specimens, and while the skull closely resembles a T. rex, it is only about half the length.} & \factcard{2}{1946}{cut.} \\ \hline
\factcard{3}{1965}{Some small tyrannosaur skeletons from Mongolia, once thought to be separate species, are young Tarbosaurus.} & \factcard{4}{1988}{The bones of the 1942 skull appear fused to one another.} \\ \hline
\factcard{5}{1988}{The 1942 skull has 15 teeth in one side of its upper jaw. Adult T. rex have 11 or 12.} & \factcard{6}{1999}{The joints between the bones of the 1942 skull are open, not fused.} \\ \hline
\factcard{7}{1999}{In Gorgosaurus, young animals have more teeth than adults: tooth number falls as the animal grows.} & \factcard{8}{2001}{Jane, a second small tyrannosaur from Hell Creek, is more than half a skeleton with a nearly complete skull.} \\ \hline
\end{tabular}
\newpage
\centering
\begin{tabular}{|@{}c@{}|@{}c@{}|}\hline
\factcard{9}{2003}{In Albertosaurus, Gorgosaurus and Daspletosaurus, young animals have low, narrow skulls and adults have deep, heavy skulls.} & \factcard{10}{2011}{A juvenile Tarbosaurus has about the same number of teeth as adult Tarbosaurus.} \\ \hline
\factcard{11}{2013}{The Dueling skeleton, a third small tyrannosaur from Hell Creek, has hand proportions and a wishbone (furcula) shape that differ from T. rex.} & \factcard{12}{2020}{Bone sections from Jane show growth rings and bone tissue typical of an animal that was still growing.} \\ \hline
\factcard{13}{2024}{In a family-tree analysis of anatomical characters, the 1942 skull and Jane may fall outside Tyrannosauridae, the family that contains T. rex.} & \factcard{14}{Jun 2025}{In tyrannosaur growth series, many characters used to build family trees differ between young and adult animals of one species.} \\ \hline
\factcard{15}{Oct 2025}{The Dueling skeleton's skull is 71 centimeters long, longer than the 1942 skull's 57.} & \factcard{16}{Oct 2025}{The outermost growth rings in the Dueling skeleton's bone are tightly packed, a pattern seen when growth has nearly stopped.} \\ \hline
\end{tabular}
\newpage
\centering
\begin{tabular}{|@{}c@{}|@{}c@{}|}\hline
\factcard{17}{Oct 2025}{Growth rings put the Dueling skeleton at 17 to 22 years old at death, and Jane at 8 to 14.} & \factcard{18}{Oct 2025}{The Dueling skeleton has 16 teeth in the left side of its upper jaw and 17 in the right. Adult T. rex have 11 or 12.} \\ \hline
\factcard{19}{Oct 2025}{The arm and hand bones of the Dueling skeleton are proportionally larger than those of T. rex.} & \factcard{20}{Oct 2025}{In living and fossil land vertebrates, limb bones do not shrink in absolute size as the animal grows.} \\ \hline
\factcard{21}{Oct 2025}{The Dueling skeleton has fewer tail vertebrae than T. rex.} & \factcard{22}{Oct 2025}{The Dueling skeleton has several air-sac openings (pleurocoels) in its second neck vertebra and in its tail vertebrae. Jane has one in the second neck vertebra and none in the tail.} \\ \hline
\factcard{23}{Dec 2025}{Bone sections from the throat bones (hyoids) of the 1942 skull show it was a mature animal.} & \factcard{24}{Jan 2026}{Growth patterns in Jane's bone do not match the growth curve built from T. rex hindlimbs of many sizes.} \\ \hline
\end{tabular}
\end{document}
